\documentclass[11pt]{article}
\usepackage[a4paper,margin=1in]{geometry}
\usepackage{amsmath,amssymb,mathtools,bm}
\usepackage{enumitem,booktabs,array}
\usepackage{tcolorbox}
\usepackage{hyperref}
\usepackage[protrusion=true,expansion=false]{microtype}
\hypersetup{colorlinks=true,linkcolor=blue,urlcolor=blue,citecolor=blue}
\setlist[itemize]{topsep=2pt,itemsep=2pt}
\setlist[enumerate]{topsep=2pt,itemsep=2pt}
\newtcolorbox{keybox}{colback=blue!3!white,colframe=blue!40!black,boxrule=0.4pt,arc=1mm}
\newtcolorbox{alertbox}{colback=red!3!white,colframe=red!40!black,boxrule=0.4pt,arc=1mm}
\newtcolorbox{answerbox}{colback=green!3!white,colframe=green!40!black,boxrule=0.4pt,arc=1mm}
\newtcolorbox{drillbox}{colback=yellow!5!white,colframe=orange!60!black,boxrule=0.4pt,arc=1mm}
\newcommand{\EV}{\mathbb{E}}
\newcommand{\Var}{\mathrm{Var}}
\newcommand{\Cov}{\mathrm{Cov}}
\newcommand{\blank}[1]{\underline{\hspace{#1}}}

\title{E05-02: Harford (2005, JFE) \\
\large ``What Drives Merger Waves?'' --- Active Recall Workbook}
\author{Empirical Corporate Finance II --- Exam Prep}
\date{\today}
\begin{document}\maketitle

\begin{keybox}
\textbf{Paper in one sentence.} Harford shows that industry-level merger waves are driven by economic, regulatory, and technological shocks \emph{combined} with sufficient capital liquidity: shocks generate the reallocation motive, while loose capital markets (low spreads, high valuations) supply the financing. Behavioral/misvaluation alone cannot explain waves --- the neoclassical plus capital-liquidity view fits the data better.

\textbf{Placement.} Canonical paper on the determinants of industry merger waves, paired with Rhodes-Kropf et al.\ (2005) which emphasizes misvaluation. Builds on Mitchell-Mulherin (1996) industry-shock evidence.
\end{keybox}

\section*{Round 1: Complete Walk-Through}

\subsection*{1.1 Research question}
Why do mergers cluster in industry-level waves, and what triggers them: real economic shocks (neoclassical) or market misvaluation (behavioral)?

\subsection*{1.2 Competing theories}
\begin{itemize}
\item \textbf{Neoclassical (Mitchell-Mulherin 1996, Jovanovic-Rousseau 2002):} industry shocks create reallocation opportunities; waves follow economic/regulatory/technological change.
\item \textbf{Misvaluation (Shleifer-Vishny 2003, Rhodes-Kropf et al.\ 2005):} overvalued stock markets let managers buy real assets with overpriced equity.
\item Harford proposes both matter, but specifically that capital-market liquidity acts as a necessary condition.
\end{itemize}

\subsection*{1.3 Data and design}
\begin{itemize}
\item Industry-level data, 1981--2000.
\item Identify merger waves: periods with abnormal industry merger activity (top-quartile merger count years concentrated in a window).
\item Shocks measured by industry-level changes in: sales, employment, R\&D, capital expenditure, profitability, net-income volatility, ROA dispersion.
\item Capital-market liquidity: C\&I spreads, GDP growth, aggregate M\&A.
\item Misvaluation: market-to-book, prior returns, analyst-forecast revisions.
\end{itemize}

\subsection*{1.4 Empirical specification}
Multinomial logit predicting whether an industry-year is part of a merger wave, on:
\[
\Pr(\text{Wave}_{it})=f(\text{Shock}_{it},\,\text{CapLiquidity}_t,\,\text{Misvaluation}_{it}).
\]

\subsection*{1.5 Headline results}
\begin{itemize}
\item Industry shocks (economic, regulatory, technological) strongly predict waves.
\item Capital-market liquidity is required: shocks without loose capital don't produce waves.
\item Misvaluation measures lose significance after controlling for shocks and capital liquidity.
\item Waves cluster in industries with large shocks in years of high aggregate liquidity (1980s deregulation, 1990s tech boom).
\end{itemize}

\subsection*{1.6 Interpretation}
\begin{itemize}
\item Supports a neoclassical-plus-capital-liquidity view: reallocation needs both a shock and cheap capital.
\item Misvaluation may matter for \emph{type} (stock vs.\ cash) of mergers but not for wave \emph{timing}.
\item Explains historical waves (railroad consolidation, conglomerate era, 1980s takeover wave, 1990s tech).
\end{itemize}

\subsection*{1.7 Caveats}
\begin{alertbox}
\begin{itemize}
\item Variables are correlated: separating shocks from misvaluation is hard.
\item Sample is US-centric and pre-dates major 2000s regulatory shifts.
\item Wave identification uses statistical cutoffs; alternative thresholds may affect inference.
\end{itemize}
\end{alertbox}

\section*{Round 2: Level 1 Fill-in}
\begin{enumerate}
\item Two competing theories: \blank{2cm} vs.\ misvaluation.
\item Capital-market \blank{2cm} required for waves to occur.
\item Shocks: economic, \blank{2cm}, technological.
\item Sample period: \blank{1cm}--\blank{1cm}.
\item Misvaluation loses \blank{2cm} after controls.
\end{enumerate}

\section*{Round 3: Level 2 Prompts}
\begin{enumerate}
\item State the two competing theories of merger waves.
\item Describe the variables and empirical specification.
\item Report the main findings on shocks, liquidity, and misvaluation.
\item Discuss the implications for the neoclassical-vs-behavioral debate.
\item Relate Harford (2005) to Rhodes-Kropf-Robinson-Viswanathan (2005).
\end{enumerate}

\section*{Round 4: Minimal Scaffolding}
From a blank page: (i) waves, (ii) theories, (iii) specification, (iv) results, (v) implications.

\section*{Round 5: Blank Page}
Essay: why do merger waves cluster in industries and time?

\vspace{6pt}
\begin{drillbox}
\textbf{Speed Drill.} (1) Two theories. (2) Necessary condition. (3) Shock types. (4) Effect of misvaluation controls. (5) Policy lesson.
\end{drillbox}

\section*{Quick Reference \& Answer Key}
\begin{answerbox}
\begin{itemize}
\item \textbf{Theories.} Neoclassical vs.\ misvaluation.
\item \textbf{Finding.} Shocks + capital liquidity drive waves.
\item \textbf{Misvaluation.} Insignificant after controls.
\item \textbf{Lesson.} Neoclassical + liquidity $>$ pure behavioral.
\end{itemize}
\end{answerbox}

\begin{keybox}
\textbf{30-second exam pitch.} Harford (2005) asks what causes industry merger waves and tests two theories: neoclassical (economic/regulatory/technological shocks drive reallocation) vs.\ misvaluation (overvalued equity lets managers cheaply buy real assets). Using industry-year data 1981--2000, with shocks measured by changes in sales, R\&D, capex, profitability, and employment, and capital-market liquidity proxied by C\&I spreads and aggregate M\&A, he runs multinomial logits predicting wave years. Shocks strongly predict waves, and capital-market liquidity is a near-necessary condition: shocks without loose capital do not produce waves. Misvaluation measures lose significance once shocks and liquidity are controlled for. The paper establishes the ``neoclassical-plus-capital-liquidity'' view of waves and contextualizes historical merger episodes (1980s deregulation, 1990s tech). It complements Rhodes-Kropf et al.\ (2005) on wave \emph{composition} while reassigning wave \emph{timing} to real shocks and liquidity.
\end{keybox}

\end{document}
